\chapter{El AST de hipótesis}
\label{ch:hypothesis}

Cada hunt es guiado por un valor \code{Hypothesis} parseado. El
AST es deliberadamente chico --- un vector plano de pasos con
listas de predicados por paso, más cláusulas opcionales de
agregación y ventana temporal adosadas a la cadena completa.

\section{Tipos}

\begin{lstlisting}[style=rust]
pub struct Hypothesis {
    pub name:           String,
    pub steps:          Vec<HypothesisStep>,
    pub k_simplicity:   usize,            // default 1
    pub aggregations:   Vec<AggClause>,   // clausulas HAVING
    pub window_seconds: Option<u64>,      // clausula WITHIN
}

pub struct HypothesisStep {
    pub origin_type:     EntityType,
    pub relation_type:   RelationType,
    pub dest_type:       EntityType,
    pub edge_predicates: Vec<Predicate>,  // sobre metadata de arista
    pub dest_predicates: Vec<Predicate>,  // sobre metadata de destino
    pub origin_binding:  Option<String>,  // `(User u)` para HAVING
    pub dest_binding:    Option<String>,
}

pub enum Predicate {
    Eq      { key: String, value: String },
    Substr  { key: String, pattern: String },
    Regex   { key: String, pattern: String },
    In      { key: String, values: Vec<String> },
}

pub struct AggClause {
    pub op:     AggOp,         // Count, CountDistinct, Min, Max, Sum, Avg
    pub column: AggColumn,     // nombre binding o indice de paso
    pub cmp:    AggComparison, // >=, <=, >, <, ==
    pub value:  f64,
}
\end{lstlisting}

\section{Invariantes}

\begin{invariant}[Cantidad de pasos $\geq 1$]
Una hipótesis debe tener al menos un paso; el parser del DSL
rechaza input vacío con \code{ApiError::InvalidInput}. Esto
libera al matcher de tratar el caso de cero aristas como una rama
especial.
\end{invariant}

\begin{invariant}[$k$-simplicidad positiva]
$k_{\text{simplicity}} \geq 1$. Cero admitiría sólo el camino
vacío y violaría el sentido de ``máximo de visitas por vértice''.
\end{invariant}

\begin{invariant}[Unicidad de bindings]
Si \code{origin\_binding} o \code{dest\_binding} es \code{Some(b)}
en algún paso, $b$ no se repite en ninguna otra posición de
binding dentro de la hipótesis. El parser lo hace cumplir en
parse time así el matcher puede resolver bindings por nombre
durante la agregación sin desambiguación.
\end{invariant}

\section{Chequeo estático de compatibilidad}

Antes de correr un hunt, el motor puede cruzar la hipótesis
contra el schema de relaciones del grafo para ver si algún paso
es estructuralmente imposible sobre el dataset cargado --- p.~ej.\
un paso \code{User -[Auth]-> Host} contra un grafo puramente de
identidad cloud donde no existen entidades \code{Host}. Esto es
\textsc{Check-Catalog-Compatibility}
(Capítulo~\ref{ch:catalog-yaml}), que devuelve un
\code{CatalogEntryWithStatus} marcando los tipos de entidad o
relación faltantes por paso.

\begin{algorithm}
\caption{\textsc{Check-Step-Compatibility}}\label{alg:check-step}
\KwIn{schema $\Sigma$ (multiconjunto de triples $(\text{source}, \text{rel},
\text{target})$), paso $s$}
\KwOut{\code{bool} compatible, conjunto de ejes faltantes}
\For{cada $(\text{src}, r, \text{dst}) \in \Sigma$}{
    \If{$(\text{src} = s.\text{origin\_type} \vee s.\text{origin\_type} = *)$
       $\wedge (r = s.\text{relation\_type} \vee s.\text{relation\_type} = *)$
       $\wedge (\text{dst} = s.\text{dest\_type} \vee s.\text{dest\_type} = *)$}{
        \Return $(\mathrm{true}, \emptyset)$\;
    }
}
$\text{faltantes} \gets $ conjunto de ejes ausentes en todo triple de $\Sigma$\;
\Return $(\mathrm{false}, \text{faltantes})$\;
\end{algorithm}

\begin{complexity}
$\Ord{|\Sigma|}$ por paso, $\Ord{|\Sigma| \cdot \ell}$ por
hipótesis, donde $\ell$ es la cantidad de pasos. $|\Sigma|$ está
acotado en la práctica por $|\mathcal{T}_V|^2 \cdot
|\mathcal{T}_E|$ (hay pocas decenas de tipos de entidad/relación),
así que el chequeo es efectivamente constante.
\end{complexity}

\begin{inthecode}
\filepath{graph\_hunter\_core/src/hypothesis.rs} --- \code{Hypothesis},
\code{HypothesisStep}, \code{Predicate}.\\
\filepath{graph\_hunter\_core/src/catalog/mod.rs:105} ---
\code{check\_catalog\_compatibility}.
\end{inthecode}
